\documentclass[11pt,a4paper]{article}

\usepackage[utf8]{inputenc}
\usepackage[T1]{fontenc}
\usepackage{amsmath,amssymb}
\usepackage{graphicx}
\usepackage{booktabs}
\usepackage[margin=2.5cm]{geometry}
\usepackage{siunitx}
\usepackage[colorlinks=true,linkcolor=blue,citecolor=blue,urlcolor=blue]{hyperref}

\newcommand{\csitl}{CsI(Tl)}
\newcommand{\lext}{\Lambda_{\rm ext}}
\newcommand{\labs}{\Lambda_{\rm abs}}
\newcommand{\lray}{\Lambda_{\rm R}}
\newcommand{\sa}{\sigma_\alpha}

\title{Calibration of the CryspLight optical model on \csitl{}, cold pure CsI and BGO:\\
the absorption/scattering split of the measured extinction}
\author{J.J.~G\'omez-Cadenas \and Claude (Anthropic)}
\date{July 13, 2026 \\[2mm] \small Version 0.3 --- calibration note (phase 1)}

\begin{document}
\maketitle

\begin{abstract}
We calibrate the optical transport model of CryspLight --- the light-propagation
simulation for the CRYSP crystals, specified in \texttt{crysp\_light\_metal.tex} --- on
\csitl{}, the best-measured scintillating crystal in the literature. Two complementary
data sets are used: the spectrophotometric bulk-extinction measurement of Knyazev et
al.~\cite{knyazev2021} ($\lext = 311.5$~mm at the 550~nm emission peak), and the
absolutely calibrated light-yield ladder of Brose, Dahlinger and
Schubert~\cite{brose1998}, who measured the \emph{same} $36$~cm crystal bare and under
successive Teflon wrappings. The central finding: the wrapping ladder is
\textbf{irreproducible} if the measured extinction is treated as pure absorption --- the
simulated wrapped/bare light-output ratio saturates near 2.2 against a measured 3.9 ---
and is reproduced, with physically consistent effective Teflon reflectivities and the
literature surface roughness, only when the extinction is split as $\sim$90\% elastic
scattering and $\sim$10\% true absorption: $\labs \simeq 3.1$~m, $\lray \simeq 0.35$~m.
Scattering redirects trapped light instead of killing it; absorption does not. The same
split resolves, at least qualitatively, the total-light-output sign discrepancy reported
by Knyazev et al.\ between their measurement and their (absorption-only) simulation.
Applied to the CRYSP geometry ($48\times48\times37.2$~mm$^3$, thick Teflon, grease-coupled
$8\times8$ SiPM array), the calibrated parameters predict that $88.8\%$ of the emitted
scintillation light reaches the sensors ($84.6$--$93.4\%$ across the split bracket
$f = 0.85$--$0.95$) --- about $2.5\times$ more than the naive reading of the measured
``absorption length'' would give.
We then calibrate \emph{cold pure CsI} on the CRYSP measurement itself
(Soleti et al.~\cite{soleti2024}, $3\times3\times20$~mm$^3$ crystals at 104~K): with the
intrinsic yield fixed at $10^5$~photons/MeV, the measured photoelectron count requires a
$\sim$10\% bulk effect, degenerate along a $(\lext, f)$ line that a single measurement
cannot break; adopting the \csitl{} analogue ($\lext = 0.3$~m, $f = 0.9$) as the default,
the CRYSP-geometry collection is $83.0\%$, with a documented 62--83\% bracket pending a
bare-crystal re-measurement or a transmittance scan.
Finally, BGO at 195~K is assembled from three sources: scintillation parameters from
Ding--Liu and Gironnet et al.\ (with a PDE-consistency argument that selects
$Y_{77\,\rm K} = 22.7$k over the published 29k), a surface/wrap validation on the
Ding--Liu 6~mm cube (reproduced with a Teflon-tape reflectivity of 0.90, consistent
with the \csitl{} thin-tape ladder), and a bulk lower bound from the Mao--Zhang--Zhu
transmittance compilation ($\lext \gtrsim 1.5$--2~m at 480~nm). The CRYSP BGO baseline
becomes $88.0\%$ (84--90\% for $\lext = 1$--3~m at $f = 0.9$; pure-absorption floor
45\%). All numbers regenerate from the \texttt{scripts/*\_calib.jl} campaigns with
their \texttt{runs/*.toml} configurations.
\end{abstract}

\tableofcontents

%=====================================================================
\section{Purpose and strategy}
\label{sec:purpose}

The CryspLight design note establishes that the predicted light collection of a wrapped
crystal is dominated by parameters that are poorly known a priori: the bulk absorption
length at the emission wavelength, the surface micro-roughness $\sa$, and --- as this
note shows --- the fraction of the measured bulk attenuation that is scattering rather
than absorption. Rather than guessing, we calibrate the full transport model on
\csitl{}: a room-temperature crystal, distinct from the cryogenic pure CsI of the CRYSP
detector (the two are kept as separate materials throughout,
\texttt{data/csitl.toml} vs \texttt{data/csi.toml}), but the one crystal for which
decades of calorimetry provide precision optical data. The strategy:

\begin{enumerate}
  \item take the \emph{bulk extinction} and the \emph{surface characterization} from
    Knyazev et al.~\cite{knyazev2021} and Janecek--Moses~\cite{janecek2010};
  \item confront the model with the \emph{wrapping ladder} of Brose et
    al.~\cite{brose1998} --- six light-yield measurements of one identical crystal under
    different wraps, which cancel the crystal's intrinsic yield, the photodetector
    response, and the reference normalization exactly;
  \item fit the minimal set of unknowns (the absorption/scattering split $f$, and one
    effective reflectivity per wrap configuration) and judge the result by
    over-constraint: the fitted reflectivities must come out physical without being
    forced to.
\end{enumerate}
Once the machinery reproduces \csitl{}, the identical procedure applies to cold pure
CsI, anchored on the CRYSP measurement itself~\cite{soleti2024}.

%=====================================================================
\section{Input data}
\label{sec:data}

\subsection{Bulk extinction and surfaces (Knyazev et al.)}

Knyazev et al.~\cite{knyazev2021} (CALIFA/R3B) measured, with a spectrophotometer over
300--750~nm on Amcrys-h \csitl{}, an ``absorption length'' rising from 28 to 34~cm
across the 500--600~nm emission band --- $311.5$~mm at the peak for their test crystal,
with a $\pm15\%$ ($1\sigma$) crystal-to-crystal spread --- together with the index chain
$n_{\rm CsI} = 1.80$, coupler 1.41, window 1.53, and ESR foil reflectivity 99.52\%.
Crucially, a collimated-beam transmission measurement counts scattered light as loss:
what is measured is the \textbf{extinction}
\begin{equation}
  \frac{1}{\lext} \;=\; \frac{1}{\labs} + \frac{1}{\lray},
  \label{eq:split}
\end{equation}
and no experiment has separated the two terms for \csitl{}. We parameterize the split by
the scattering fraction $f$: $\labs = \lext/(1-f)$, $\lray = \lext/f$. Knyazev et al.\
also characterized their surfaces by AFM (slope RMS 0.15 polished, 0.54 lapped at 100~nm
sampling) and, importantly for us, reported that their absorption-only simulation
\emph{reproduces the uniformity trend but gets the sign of the total-light-output change
under lapping wrong}, naming unmodelled scattering (``crystal dimming'') as the suspected
cause. The optically effective roughness values are those of
Janecek--Moses~\cite{janecek2010}: $\sa = 1.3^\circ$ for mechanically polished,
$\sim$$18^\circ$ for ground surfaces.

\subsection{The absolutely calibrated wrapping ladder (Brose et al.)}
\label{sec:brose}

Brose et al.~\cite{brose1998} (TU Dresden, early BABAR R\&D) measured Kharkov \csitl{}
crystals --- tapered, $6\times5$~cm$^2$ front, $6\times6$~cm$^2$ rear, 36~cm long
(19.4~$X_0$) --- with two absolutely linked scales: a PMT scale LY (Hamamatsu R669,
minimal air gap, 5~cm diameter photocathode; LY quoted relative to a 25~mm reference
cube, averaged over an 18-position $^{137}$Cs scan along the crystal), and a photodiode
scale PY in photoelectrons per MeV, calibrated by absorbing 59.5~keV photons directly in
the diode (16480 e--h pairs; accuracy $\pm100\;{\rm e_0/MeV}$; QE $85\%$ at 560~nm). The
bridge: ${\rm LY} = 74\% \leftrightarrow {\rm PY} = 7600\;{\rm e_0/MeV}$. The
measurements used here (Fig.~7 of~\cite{brose1998}, one and the same polished crystal):

\begin{table}[htb]
  \centering
  \caption{The Brose et al.\ wrapping ladder and the ratios used as fit targets.}
  \label{tab:ladder}
  \begin{tabular}{lcc}
    \toprule
    Configuration & LY & LY/LY$_{\rm bare}$ \\
    \midrule
    bare (internal reflection only)        & 15\% & 1    \\
    1 layer Teflon, 38~$\mu$m              & 34\% & 2.27 \\
    3 layers Teflon                        & 53\% & 3.53 \\
    4 layers Teflon ($\sim$160~$\mu$m)     & 59\% & 3.93 \\
    3 layers + 5~$\mu$m Al foil            & 59\% & 3.93 \\
    1 layer Teflon, 200~$\mu$m (monolayer) & 41\% & 2.73 \\
    \bottomrule
  \end{tabular}
\end{table}

Also used qualitatively: an air gap at the photodiode costs $\sim$15\% of the
photoelectron yield versus glued coupling; removing the reflector that covers the
uninstrumented rear-face area costs $\sim$30\%. Both are direct probes of physics the
model must contain (the Fresnel escape cone at the readout; the recycling of photons by
the readout-face surround).

%=====================================================================
\section{Simulation setup (phase 1)}
\label{sec:setup}

The Brose PMT configuration is modelled in CryspLight as follows; every number below is
in \texttt{runs/csitl\_calib.toml}.

\begin{itemize}
  \item \textbf{Geometry}: straight-box approximation of the tapered crystal,
    $60\times55\times360$~mm$^3$ (mean cross-section). The $0.8^\circ$ taper --- which
    under specular TIR walks photon angles deterministically, an unlocking mechanism
    analogous to surface roughness --- is deferred to phase~2; its expected effect is
    noted in Sec.~\ref{sec:results}.
  \item \textbf{Readout}: a disc of radius 25~mm (the photocathode) on the
    $6\times6$~cm$^2$ rear face, coupled through an \emph{air gap}
    ($n_{\rm coupling}=1$; escape cone $\arcsin(1/1.79) = 34.0^\circ$), PDE $= 1$
    (all targets are relative). Photons transmitted outside the disc hit the
    \emph{surround}: absorbing in the bare configuration (nothing is there), and ---
    an essential detail --- \textbf{Teflon in the wrapped configurations}, since the
    whole crystal is wrapped except the PMT window; the surround is then a Lambertian
    reflector with the same $R$ as the wrap.
  \item \textbf{Source}: nine equal-weight isotropic point sources along the crystal
    axis ($z = 20$--$340$~mm), approximating the 18-position scan average; delta
    emission time. $4\times10^5$ photons per configuration.
  \item \textbf{Surfaces}: UNIFIED back-painted (air gap) with the specular lobe;
    $\sa \in \{1.3^\circ, 8.5^\circ\}$ (optical polish vs the AFM upper bracket).
  \item \textbf{Bulk}: $\lext = 311.5$~mm split by Eq.~\ref{eq:split} with
    $f \in \{0, 0.3, 0.5\}$ (first pass) and $\{0.8, 0.9, 0.95\}$ (second pass);
    Rayleigh phase function $\propto 1+\cos^2\theta$.
  \item \textbf{Fit}: for each $(f, \sa)$, simulate the bare crystal; for each wrap
    configuration, scan the wrap reflectivity $R$ and interpolate the effective
    $R_{\rm eff}$ at which the simulated ratio to bare equals the measured one
    (Table~\ref{tab:ladder}). The wrap model is Lambertian behind the air gap
    (\textsc{unified} \texttt{groundbackpainted}).
\end{itemize}

The transport engine itself is validated independently by the analytic anchors of the
design note (exact photon conservation; the idealized-specular escape-cone limit
$1-\cos\theta_c'$; Rayleigh-restored ergodicity; Fresnel closed forms; sampling
distributions), enforced in the package test suite.

%=====================================================================
\section{Results}
\label{sec:results}

\subsection{Pure absorption is excluded}

With $f \le 0.5$, the simulated wrapped/bare ratio saturates near $2.2$ --- below even
the single-thin-layer target of 2.27, and far from the 4-layer target of 3.93 ---
\emph{regardless of the wrap reflectivity} (the ratio is nearly flat in $R$ up to
$R = 0.99$). The mechanism: light recycled by the wrap random-walks and must then
survive the bulk; with $\labs \simeq \lext = 31$~cm in a 36~cm crystal it does not, so
the wrap cannot matter, contrary to observation. The same conclusion follows from the
measurement alone: Brose's \emph{bare} crystal shows only a few-percent position
dependence of the light output along 36~cm, incompatible with 31~cm of true absorption.

\subsection{The scattering-dominated split}

\begin{table}[htb]
  \centering
  \caption{Fitted effective Teflon reflectivities per grid point
  ($4\times10^5$ photons per run; statistical precision on the ratios
  $\lesssim 1\%$). ``$>$'' / ``$<$'': the target ratio is not reached within the
  scanned $R$ range (above its ceiling / below its floor). The preferred point is
  \textbf{bold}.}
  \label{tab:grid}
  \begin{tabular}{ccccccc}
    \toprule
    $f$ & $\sa$ & bare frac. & $R_{\rm eff}$(1$\times$38$\mu$m) &
    $R_{\rm eff}$(3$\times$) & $R_{\rm eff}$(4$\times$) & $R_{\rm eff}$(200$\mu$m mono) \\
    \midrule
    0.80 & $1.3^\circ$ & 0.054 & 0.867 & $>0.995$ & $>0.995$ & 0.966 \\
    0.80 & $8.5^\circ$ & 0.054 & 0.778 & 0.963    & $>0.995$ & 0.868 \\
    \textbf{0.90} & $\mathbf{1.3^\circ}$ & \textbf{0.060} & \textbf{0.754} &
      \textbf{0.934} & \textbf{0.965} & \textbf{0.840} \\
    0.90 & $8.5^\circ$ & 0.058 & $<0.75$ & 0.880 & 0.910 & 0.793 \\
    0.95 & $1.3^\circ$ & 0.064 & $<0.75$ & 0.866 & 0.896 & 0.781 \\
    0.95 & $8.5^\circ$ & 0.061 & $<0.75$ & $<0.85$ & $<0.88$ & $<0.78$ \\
    \bottomrule
  \end{tabular}
\end{table}

Table~\ref{tab:grid} shows the second-pass grid. The point $f = 0.90$,
$\sa = 1.3^\circ$ is singled out on three independent grounds:
\begin{enumerate}
  \item it is the only point at which \emph{all four} ladder rungs are matched within
    the physical reflectivity range;
  \item the fitted $R_{\rm eff}$ come out \textbf{monotone in Teflon thickness, below
    saturation, and with the 200~$\mu$m monolayer correctly ranked between one and
    three thin layers and below the four-layer multilayer} --- the measured ordering,
    which the fit was free to violate;
  \item the preferred $\sa$ is exactly the Janecek--Moses polished value, which was
    \emph{not} fitted.
\end{enumerate}
Neighbouring points fail in opposite directions ($f = 0.8$: the thick rungs need
$R > 0.995$; $f = 0.95$: even the scan floor overshoots the thin rungs), bracketing
$f \approx 0.85$--$0.92$ at this fidelity. The remaining tension --- the simulated bare
absolute collection of 6.0\% against the $\sim$8--10\% implied by ${\rm LY} = 15\%$
under reasonable assumptions for the reference cube --- has the sign and rough size
expected from the missing taper, and is the primary target of phase~2.

\subsection{Final calibrated parameters}

\begin{table}[htb]
  \centering
  \caption{The calibrated \csitl{} optical model (550~nm, room temperature) ---
  the content of \texttt{data/csitl.toml}.}
  \label{tab:final}
  \small
  \begin{tabular}{lll}
    \toprule
    Parameter & Value & Origin \\
    \midrule
    Refractive index $n$ & 1.79 & \cite{brose1998,knyazev2021} \\
    Extinction $\lext$ & 311.5~mm ($\pm15\%$ batch) & measured \cite{knyazev2021} \\
    Scattering fraction $f$ & $0.90$ (bracket 0.85--0.92) & this fit \\
    $\Rightarrow$ absorption $\labs$ & $3115$~mm & $\lext/(1-f)$ \\
    $\Rightarrow$ scattering $\lray$ & $346$~mm & $\lext/f$ \\
    Surface $\sa$ (polished, air gap) & $1.3^\circ$ & \cite{janecek2010}, confirmed \\
    Teflon $R_{\rm eff}$ (1 / mono / 3 / 4 layers) &
      0.754 / 0.840 / 0.934 / 0.965 & this fit \\
    Teflon saturated ($\gtrsim 0.5$--1~mm) & 0.99 & \cite{janecek2008} \\
    Light yield / decay & $54\,000$~ph/MeV; $\sim$730~ns (64\%) + 3.34~$\mu$s &
      literature \\
    \bottomrule
  \end{tabular}
\end{table}

The physics statement: \textbf{the measured ``absorption length'' of \csitl{} is
extinction, and about nine parts in ten of it is elastic scattering}; true absorption is
at the metre scale. This simultaneously accounts for the Brose wrapping gains, the
near-flat bare-crystal uniformity, the long ``attenuation lengths'' historically
extracted from calorimeter uniformity fits, and the direction of the
simulation-vs-measurement discrepancy Knyazev et al.\ left open.

%=====================================================================
\section{Application to the CRYSP geometry}
\label{sec:crysp}

With Table~\ref{tab:final} and the standard CRYSP test of the design note --- a
$48\times48\times37.2$~mm$^3$ crystal, thick (saturated) Teflon on five faces,
$\sa = 1.3^\circ$, grease coupling ($n_g = 1.45$), fully active $8\times8$ SiPM face,
isotropic point source at the centre, PDE $= 1$:
\begin{equation}
  \boxed{\;\varepsilon_{\rm coll} = 88.8\% \quad (84.6\text{--}93.4\% \text{ for }
  f = 0.85\text{--}0.95)\;}
\end{equation}
with 10.0\% absorbed in the bulk and 1.2\% at the wrap; detected photons average
$\sim$11 wall interactions. The bracket is narrow because at 37~mm the crystal is short
compared with \emph{both} lengths of the split. In photoelectrons: a 511~keV photopeak
emits $\sim$27\,600 photons, $\sim$24\,500 reach the sensors, and a realistic
${\rm PDE} \approx 0.4$ at 550~nm gives $\sim$$10^4$ photoelectrons per event. Had the
measured 31~cm been used as true absorption, the same test returns roughly a third of
the light: the split is worth a factor $\sim$2.5 in predicted signal.

For the CRYSP detector proper the open question transfers to \emph{cold pure CsI at
340--350~nm} (a separate material in the framework): if its extinction is
scattering-dominated like its doped sibling --- plausible, though the UV absorption edge
is closer --- its collection also sits near the top of this range. That split should be
calibrated on the CRYSP measurements themselves~\cite{soleti2024}, by exactly the
procedure of this note; a bare-vs-wrapped ladder on a CRYSP crystal would be the single
most informative measurement.

%=====================================================================
\section{Cold pure CsI: calibration on the CRYSP measurement}
\label{sec:coldcsi}

Cold pure CsI --- the actual CRYSP detector material, kept separate from \csitl{}
throughout (\texttt{data/csi.toml}) --- has no published bulk-optics data at its
340--350~nm cryogenic emission. Its calibration anchor is the CRYSP measurement
itself~\cite{soleti2024}: two $3\times3\times20$~mm$^3$ AMCRYS crystals (1.1~$X_0$,
aspect ratio $\sim$7, polished), wrapped in \emph{three layers of 80~$\mu$m PTFE tape}
(unsaturated: $R \approx 0.93$--0.97), each read by a single $3\times3$~mm$^2$
S13360-3025CS covering the full end face in \emph{dry contact} (air gap), at $T=104$~K
with a $^{22}$Na source. Measured: $N_{\rm pe} = 4596$ at 511~keV (crystal~2:
$\sim$4383), energy resolution 6.30\% FWHM, $\tau \simeq 810$~ns.

\subsection{Anchor and setup}

Taking the intrinsic yield as fixed at $Y = 10^5$~photons/MeV (robust across the
literature), the measured collection efficiency is
\begin{equation}
  \varepsilon_{\rm det} = \frac{N_{\rm pe}}{Y \cdot E \cdot {\rm PDE}}
  = \frac{4596}{10^5 \cdot 0.511 \cdot 0.17} = 0.529
  \qquad (\text{crystal 2: } 0.505),
\end{equation}
with the cryo-spectrum-weighted ${\rm PDE} = 17\%$ as the dominant systematic (the
cryogenic PDE of this MPPC is not directly measured). The simulation
(\texttt{runs/soleti\_calib.toml}): the $3\times3\times20$ box, full-face readout through
an air gap ($n_{\rm coupling}=1$, escape cone $30.9^\circ$), five faces PTFE
(back-painted Lambertian, $\sa = 1.3^\circ$), 511~keV deposits distributed exponentially
along the axis (attenuation 23~mm, entering the face opposite the sensor), ${\rm PDE}=1$.
Two consistency checks come out well: with a near-transparent bulk our model gives
$\varepsilon = 0.56$--0.60, in agreement with the Geant4 LUT-model estimate of 0.58
quoted in~\cite{soleti2024} (two independent transport codes, same answer); and
$Y = 10^5$ then \emph{requires} a $\sim$10\% bulk effect to reach the measured 0.529.

\subsection{A one-measurement degeneracy}

Scanning $(\lext, f, R_{\rm wrap})$, the measured $\varepsilon$ is reproduced along a
degeneracy line; three representative solutions:
\begin{center}
\begin{tabular}{cccc}
  \toprule
  $\lext$ & $f$ & $R_{\rm wrap}$ & CRYSP-geometry collection \\
  \midrule
  0.3~m & 0.9 & 0.93 & 82.9\% \\
  1~m   & 0.5 & 0.95--0.97 & 63.9\% \\
  3~m   & 0   & 0.93 & 62.2\% \\
  \bottomrule
\end{tabular}
\end{center}
The last column is the point: the degeneracy is harmless in the small crystal but
propagates to a 20-point spread in the CRYSP geometry (scattering rescues trapped light;
absorption does not). We attempted to break it with the published energy resolution via
the depth dependence of $\varepsilon$: the result is a null with its own value ---
$\varepsilon(z)$ is flat to $<0.5\%$ FWHM for \emph{all three} solutions (the diffuse
wrap fully homogenizes the light in this geometry), so the resolution budget cannot
discriminate, and conversely the uniform-collection assumption used in the intrinsic
resolution extraction of~\cite{soleti2024} is validated.

\subsection{Adopted default and how to break the degeneracy}

We adopt $\boxed{\lext = 0.3~\text{m},\; f = 0.9}$
($\labs = 3$~m, $\lray = 0.33$~m) as the default --- not as the simplest option but as
the \emph{coherent} one: the same producer family (AMCRYS/Kharkov) shows
$\lext \approx 0.31$~m at 550~nm in \csitl{} with $f \approx 0.9$, and a scattering
length flat between 340 and 550~nm is the signature of defect/inclusion (Mie-type)
scattering rather than molecular Rayleigh ($\propto \lambda^{-4}$); with no dopant band,
the metre-scale absorption is natural in the pure crystal. One physical picture, one
parameter set, both materials. The resulting CRYSP baseline is
\textbf{83.0\% collection} ($\sim$42\,000 photons at the sensor plane per 511~keV
photopeak; $\sim$11--13k photoelectrons at ${\rm PDE} = 0.25$--0.30), with the
\textbf{62--83\% bracket} kept in force until the degeneracy is broken. Ranked breakers,
the first two in-house:
\begin{enumerate}
  \item \textbf{re-measure the same $3\times3\times20$ crystals bare} (no PTFE) --- the
    Brose trick: the bare/wrapped ratio separates wall escape from bulk physics;
  \item \textbf{a spectrophotometer transmittance scan} of a pure CsI sample at
    340--350~nm (Knyazev-style): $\lext$ measured directly pins $f$ from the existing
    wrapped point;
  \item a global fit adding published cryo-CsI setups of different aspect
    ratios~\cite{ding2022,wang2024};
  \item the forthcoming CRYSP SiPM matrices / monolithic crystals: transverse light-map
    shapes are strongly $f$-sensitive.
\end{enumerate}

%=====================================================================
\section{BGO at 195~K: parameters, validation, and the CRYSP prediction}
\label{sec:bgo}

BGO (the CRYSP alternative crystal, operated at dry-ice temperature) is assembled from
three literature sources plus one validation campaign; materials
\texttt{data/bgo.toml} (195~K) and \texttt{data/bgo\_77k.toml}.

\subsection{Scintillation parameters}

Two independent chains agree on $Y(195~{\rm K}) = 14$k~ph/MeV: the Ding--Liu
measurement~\cite{ding2024} ($5.2\pm0.3$~PE/keV system yield, converted with their
simulated collection and PDE) and the Gironnet quenching curve~\cite{gironnet2008}
anchored on Moszy\'nski's room-temperature 6.9k. The decay time at exactly 195~K is
$\tau = 1768$~ns (Ding--Liu waveform fit; Piltingsrud and Gironnet bracket it from
above at 200~K). At 77~K, a consistency argument selects the yield: Ding--Liu derived
29k assuming temperature-independent PDE at fixed bias, but their own single-PE
calibration shows the gain (hence overvoltage, hence PDE) rising by $\times1.24$ from
195 to 77~K --- and Gironnet's $22.7\pm2$k times 1.24 reproduces their measured
system-yield ratio (10.5/5.2 = 2.02) exactly. We adopt $Y(77~{\rm K}) = 22.7$k with the
two-component decay 1.4~$\mu$s (24\%) + 8.7~$\mu$s (76\%). Unlike CsI, the emission
stays at the Bi$^{3+}$ 480~nm band at all temperatures.

\subsection{Surface/wrap validation on the Ding--Liu cube}

The Ding--Liu geometry --- a $6\times6\times6$~mm$^3$ cube, \emph{unpolished} (rough)
surfaces, four sides in multiple layers of Teflon tape, two SiPMs dry-coupled to the
remaining faces --- is simulated with the two-sided readout
(\texttt{runs/ding\_calib.toml}; deposits exponential, 0.6~mm, inside the entrance side
face). With $Y = 14$k and ${\rm PDE} = 0.45$, the measured collection is
$\varepsilon = 0.825 \pm 0.08$. The campaign confirms what the geometry implies: at
6~mm, $\sa$ (12$^\circ$ vs 18$^\circ$) and $\lext$ (1.5 vs 2~m) move $\varepsilon$ by
$<0.5\%$ --- the cube measures the \emph{tape}. The measured value is reproduced at
$R_{\rm tape} = 0.90$ (their Geant4 MC value of 0.80 corresponds to $R = 0.88$),
squarely inside the \csitl{} thin-tape ladder of Sec.~\ref{sec:results}
(0.84 for a 200~$\mu$m monolayer, 0.93 for three 38~$\mu$m layers): the model
reproduces the measurement with no new physics. This is a validation of the
rough-surface + tape + air-gap branch, not a bulk constraint.

\subsection{Bulk and the CRYSP prediction}

Mao, Zhang and Zhu~\cite{mao2008} measured the transmittance of a 1.68~cm BGO cube at
the Fresnel-limit level, giving $\lext \gtrsim 1.5$--2~m at 480~nm (BGO's emission sits
deep in its transparency region; cutoff 315~nm). We adopt $\lext = 2$~m with $f = 0.9$
(the community default --- extinction $\approx$ scattering at the emission wavelength
of a clean crystal), noting that \emph{for BGO the split is unmeasured}: no
wrapping-ladder or dedicated-geometry data exist to break it, and the same in-house
breakers as for cold CsI apply. In the CRYSP geometry
($48\times48\times22.4$~mm$^3$, thick Teflon, grease-coupled $8\times8$ array,
$\sa = 1.3^\circ$, ${\rm PDE}=1$):
\begin{center}
\begin{tabular}{lc}
  \toprule
  bulk hypothesis & CRYSP collection \\
  \midrule
  $\lext = 2$~m, $f = 0.9$ (\textbf{adopted}) & \textbf{88.0\%} \\
  $\lext = 1$--3~m, $f = 0.9$ & 84.5--90.0\% \\
  $\lext = 2$~m, $f = 0$ (pure absorption) & 45.3\% \\
  $\lext = 1$~m, $f = 0$ (old placeholder) & 37.4\% \\
  \bottomrule
\end{tabular}
\end{center}
The tracked baselines become 88.0\% (Teflon) and 70.0\% (ESR --- legitimate for BGO's
480~nm emission; the specular wrap pays $\sim$18\% wall losses against the diffuse
wrap's 1\%, a first quantitative argument for Teflon over ESR in the CRYSP BGO
option). In photoelectrons: $14\,000 \times 0.511 \times 0.88 \times {\rm PDE}(0.45)
\approx 2\,800$ per 511~keV photopeak at 195~K.

%=====================================================================
\section{Caveats and next steps}
\label{sec:caveats}

\begin{enumerate}
  \item \textbf{Taper (phase 2).} The straight-box approximation is the leading known
    systematic; it plausibly accounts for the bare-crystal 6\% vs 8--10\% gap and will
    sharpen the $f$ bracket. Requires generalizing the wall intersection and surface
    frames to non-axis-aligned planes.
  \item \textbf{Source placement.} On-axis points, whereas the measured deposits sit a
    few cm below the top surface; a second-order effect on the ladder ratios, free to
    fix via the config.
  \item \textbf{Scalar optics.} Single wavelength at the emission peak; the
    28--34~cm band variation and the PDE$(\lambda)$ weighting enter only at the few-\%
    level for these ratios.
  \item \textbf{Crystal provenance.} Brose (1990s Kharkov) and Knyazev (Amcrys-h)
    crystals differ; the $\pm15\%$ batch spread on $\lext$ is folded into the $f$
    bracket only informally.
  \item \textbf{Monolayer/multilayer.} The ordering is \emph{reproduced through fitted}
    $R_{\rm eff}$, not predicted ab initio (inter-layer air gaps are not modelled
    microscopically).
  \item \textbf{Phase 3.} The absolute photoelectron scale (the photodiode geometry:
    glued coupling, Tyvek-covered surround, PY = 7600~e$_0$/MeV) as an end-to-end
    absolute closure test, after the taper.
\end{enumerate}

\section*{Reproducibility}

\begin{itemize}
  \item \texttt{scripts/csitl\_calib.jl} + \texttt{runs/csitl\_calib.toml} --- the
    \csitl{} campaign (grid, cases, targets);
  \item \texttt{scripts/soleti\_calib.jl} + \texttt{runs/soleti\_calib.toml} --- the
    cold-CsI campaign (Sec.~\ref{sec:coldcsi});
  \item \texttt{scripts/ding\_calib.jl} + \texttt{runs/ding\_calib.toml} --- the BGO
    cube campaign (Sec.~\ref{sec:bgo});
  \item \texttt{data/csitl.toml}, \texttt{data/csi.toml}, \texttt{data/bgo.toml},
    \texttt{data/bgo\_77k.toml} --- the calibrated parameters, provenance in comments;
  \item \texttt{runs/point\_csitl\_teflon.toml}, \texttt{runs/point\_csi\_teflon.toml}
    --- the CRYSP-geometry applications;
  \item \texttt{output/csitl\_calib/}, \texttt{output/soleti\_calib/} --- grid results;
  \item the package test suite (\texttt{Pkg.test}) enforces the transport anchors.
\end{itemize}

%=====================================================================
\begin{thebibliography}{9}

\bibitem{knyazev2021}
A.~Knyazev, J.~Park, P.~Golubev \emph{et al.},
``Simulations of light collection in long tapered CsI(Tl) scintillators using real
crystal surface data and comparisons to measurement,''
Nucl.\ Instrum.\ Meth.\ A \textbf{1003} (2021) 165302.

\bibitem{brose1998}
J.~Brose, G.~Dahlinger and K.~R.~Schubert,
``Properties of CsI(Tl) crystals and their optimization for calorimetry of high energy
photons,''
Nucl.\ Instrum.\ Meth.\ A \textbf{417} (1998) 311, arXiv:physics/9805010.

\bibitem{janecek2010}
M.~Janecek and W.~W.~Moses,
``Simulating scintillator light collection using measured optical reflectance,''
IEEE Trans.\ Nucl.\ Sci.\ \textbf{57} (2010) 964.

\bibitem{janecek2008}
M.~Janecek and W.~W.~Moses,
``Optical reflectance measurements for commonly used reflectors,''
IEEE Trans.\ Nucl.\ Sci.\ \textbf{55} (2008) 2432.

\bibitem{soleti2024}
S.~R.~Soleti \emph{et al.} (CRYSP),
``Cryogenic cesium iodide as a potential PET material,''
arXiv:2406.13598 (2024).

\bibitem{ding2022}
K.~Ding, J.~Liu, Y.~Yang and D.~Chernyak,
``First operation of undoped CsI directly coupled with SiPMs at 77~K,''
Eur.\ Phys.\ J.\ C \textbf{82} (2022) 344.

\bibitem{wang2024}
L.~Wang \emph{et al.},
``Reactor neutrino physics potentials of cryogenic pure-CsI crystal,''
Eur.\ Phys.\ J.\ C \textbf{84} (2024) 440.

\bibitem{ding2024}
K.~Ding and J.~Liu,
``Operation of BGO with SiPM readout at dry ice and liquid nitrogen temperatures,''
arXiv:2403.09501 (2024).

\bibitem{gironnet2008}
J.~Gironnet, V.~B.~Mikhailik, H.~Kraus, P.~de~Marcillac and N.~Coron,
``Scintillation studies of Bi$_4$Ge$_3$O$_{12}$ (BGO) down to a temperature of 6~K,''
Nucl.\ Instrum.\ Meth.\ A \textbf{594} (2008) 358.

\bibitem{mao2008}
R.~Mao, L.~Zhang and R.-Y.~Zhu,
``Optical and scintillation properties of inorganic scintillators in high energy
physics,''
IEEE Trans.\ Nucl.\ Sci.\ \textbf{55} (2008) 2425.

\end{thebibliography}

\end{document}
